\section{Vary the implementation and explain the evidence}\label{blit-practice}
In free-writing mode, keep the function interface and the supplied caller but replace the guided body. One valid approach examines each source cell, computes its destination coordinate, and bounds-checks that coordinate before a copy. Another clips once and walks only the intersection. Compare them on a fully visible sprite, a heavily clipped sprite, and a fully invisible sprite. Re-run the behavioral checks after every change.

Use the native instruction trace to count executed instructions from function entry to return. Use memory-access traces to count source reads separately from destination writes. These quantities describe simulator work for a particular input; they are not elapsed-time measurements. Report transparency patterns and placements with each measurement. A lower count for an invisible sprite does not establish a faster implementation for every visible sprite.

Trace the wrong indices caused by width/stride confusion and an incorrect clipped source origin. Design a boundary case for each. Explain why counting changed values fails for an opaque color already at its destination. A final image can conceal write-and-restore errors; inspect write history for transparency and padding. Wider accesses remain acceptable when every byte satisfies the contract.
